\documentclass{ximera}
\title{Practice Quiz 1, Approximating Scaling Factors Numerically}


\newcommand{\pskip}{\vskip 0.1 in}

\begin{document}
\begin{abstract}
Practice Quiz 1.
\end{abstract}
\maketitle


\pskip

{\bf Directions:}

\begin{itemize}
\item Give thorough explanations in complete sentences to each of the following questions. Include pictures for Questions 1 and 2 to help with your explanation. 

\item You will need a scientific calculator, but do not use any other technology.

\end{itemize}


\pskip

\begin{question} \label{QKdf3rt3cg}
Use precalculus and arithmetic, but no calculus, to give numerical evidence that suggests an answer to the question below:

\pskip

In the late afternoon, a $100$-foot tall building casts a $300$-foot long shadow on level ground. Is there a scaling factor that converts (by multiplication) a small change in the inclination angle of the setting sun to an accurate approximation of the corresponding change in the shadow's length? 

\pskip

\begin{enumerate}

\item Start by writing a complete sentence that defines a function expressing the length of the shadow in terms of the inclination angle. Include units in your definition. Use meaningful variables, \emph{not} $x$ and $y$. Then find an expression for the function.


\item If your evidence suggests that such a scaling factor exits, explain why and approximate the scaling factor to the nearest hundredth. Include units. Otherwise, if your evidence suggests that no such scaling factor exists, explain why.

\end{enumerate}
\end{question}




\begin{question} \label{QK3ef34n3}

Use precalculus and arithmetic, but no calculus, to give numerical evidence that suggests an answer to the question below:

\pskip

A rocket orbits a planet of radius $2000$km at an altitude of $400$ km above the surface. Is there a scaling factor that converts (by multiplication) a small change in orbit's radius to an accurate approximation of the corresponding change in radius of the portion of the planet's surface visible to an astronaut in the rocket? The radius is measured along the planet's surface as shown below.


\begin{onlineOnly}
    \begin{center}
\desmosThreeD{paqpozwuxm}{900}{600}
\end{center}
\end{onlineOnly}

Desmos activity available at
\href{https://www.desmos.com/3d/paqpozwuxm}{151: Spaceship above earth}
 

\pskip

\begin{enumerate}

\item Start by writing a complete sentence that defines the appropriate function. Include units in your definition. Use meaningful variables, \emph{not} $x$ and $y$. Then find an expression for the function.

\item If your evidence suggests that such a scaling factor exits, explain why and approximate the scaling factor to the nearest hundredth. Include units. Otherwise, if your evidence suggests that no such scaling factor exists, explain why.

\end{enumerate}
\end{question}


\begin{question}  \label{Qdr3t3t3m}
Use precalculus and arithmetic, but no calculus, to give numerical evidence that suggests an answer to the question below:

\pskip

The atmospheric pressure is $3$ atm at an altitude of $5$ km above the surface of the planet Krypton. The pressure is $2$ atm at an altitude of $12$km. Assume the pressure is an exponential function of altitude.

A weather balloon is $5$km above the surface. Is there a scaling factor that converts (by multiplication) a small change in the balloon's altitude to an accurate approximation of the corresponding change in atmoshperic pressure?

\pskip

\begin{enumerate}

\item Start by writing a complete sentence that defines the appropriate function. Include units in your definition. Use meaningful variables, \emph{not} $x$ and $y$. Then find an expression for the function. Include a graph of the function. 


\item If your evidence suggests that such a scaling factor exits, explain why and approximate the scaling factor to the nearest hundredth. Include units. Otherwise, if your evidence suggests that no such scaling factor exists, explain why.

\end{enumerate}

\end{question}





\end{document}